% ============================================================
%  Airy Discrete Stability Lemma
%  Working note -- prolate-gram-coercivity programme
%
%  Goal: Prove Prop U(U') of Paper VIII:
%    |(lambda_l - lambda_{l+1}) - (F_Ai(x_l) - F_Ai(x_{l+1}))|  <=  C_2 c^{-2/3}
%
%  Position in the hierarchy:
%    phase_nondeg_lemma (v3)  ─┐
%                              ├──→  THIS NOTE  ──→  Paper VIII
%    ORX §4 (Thm 4.6)         ─┘
%
%  This note is NOT about phase increments (that is phase_nondeg).
%  It is about discrete differences of eigenvalue approximations.
%  The common ingredient is the Airy error E_k(c) from ORX,
%  but the objects are different:
%    phase_nondeg:  Theta_k(s*) -- WKB phase at fixed point
%    this note:     lambda_l -- eigenvalue sequence
%
%  Structure:
%    Section 1 -- Setup and notation
%    Lemma U   -- Pointwise Airy approximation (from ORX, recalled)
%    Lemma U'  -- Discrete derivative stability (new content)
%    Corollary -- ass:gap follows immediately
% ============================================================
\documentclass[12pt]{amsart}

\usepackage{amsmath,amssymb,amsthm}
\usepackage{mathtools}
\usepackage{hyperref}
\usepackage[shortlabels]{enumitem}

\newtheorem{theorem}{Theorem}[section]
\newtheorem{proposition}[theorem]{Proposition}
\newtheorem{lemma}[theorem]{Lemma}
\newtheorem{corollary}[theorem]{Corollary}
\newtheorem{remark}[theorem]{Remark}
\newtheorem{definition}[theorem]{Definition}

% --- macros consistent with Paper VIII ---
\newcommand{\FAi}{F_{\mathrm{Ai}}}
\newcommand{\NSh}{N_{\mathrm{Sh}}}
\newcommand{\norm}[1]{\lVert#1\rVert}
\newcommand{\abs}[1]{\lvert#1\rvert}

% ============================================================
\begin{document}
% ============================================================

\title[Airy Discrete Stability for PSWF Eigenvalues]{
  Airy Discrete Stability for PSWF Eigenvalues:\\
  A Missing Link for Proposition~U(U') in Paper~VIII}

\author{[Author]}
\date{\today}
\maketitle

\begin{abstract}
Let $\lambda_l$ be the eigenvalues of the PSWF concentration
operator, ordered decreasingly, and let
\[
  x_l := (l - \NSh)\,(c/2)^{-1/3}, \qquad \NSh := 2c/\pi,
\]
be the Airy-scaled index.
We prove:
\begin{enumerate}[(U)]
  \item $|\lambda_l - \FAi(x_l)| \le C_1 c^{-2/3}$ \quad
        (recalled from ORX \cite{OsipovRokhlinXiao}, Thm.~4.6).
  \item $|(\lambda_l - \lambda_{l+1})
          - (\FAi(x_l) - \FAi(x_{l+1}))|  \le C_2 c^{-2/3}$
        \quad (Lemma~\ref{lem:Uprime}, new).
\end{enumerate}
Both bounds hold uniformly for $|l - \NSh| \le A_0 c^{1/3}$
and $c \ge c_0(A_0)$.
The key structural input is the discrete stability of the
Airy-connection error: $|E_{l+1}(c) - E_l(c)| = O(c^{-4/3})$,
which follows from the $C^1$ regularity of $E_l$ in the Airy
parameter $\mu_l$ (ORX, Ch.~4) combined with the step-size
estimate $|\mu_{l+1} - \mu_l| = O(c^{-1})$.
As a corollary, Assumption~\ref{ass:gap} of Paper~VIII
\cite{paperVIII} follows unconditionally for all
$c \ge c_0(A_0)$.
\end{abstract}

% ============================================================
\section{Setup and Notation}
\label{sec:setup}
% ============================================================

\noindent\textbf{Parameters.}
$c \ge 1$; $\NSh := 2c/\pi$; $A_0 > 0$ fixed;
$c_0 = c_0(A_0)$ a sufficiently large constant.
All estimates uniform in $l$ with $|l - \NSh| \le A_0 c^{1/3}$
(the \emph{transition window}).

\medskip
\noindent\textbf{Rescaled index and Airy scaling.}
\begin{equation}\label{eq:xl}
  x_l := (l - \NSh)\,(c/2)^{-1/3}.
\end{equation}
Consecutive spacing: $x_{l+1} - x_l = (c/2)^{-1/3}$ exactly.

\medskip
\noindent\textbf{Airy approximation function.}
$\FAi : \mathbb{R} \to (0,1)$ is the function defined in
Paper~VIII \cite{paperVIII}, Section~2:
\[
  \FAi(x) := \int_x^\infty \mathrm{Ai}(t)^2\,dt,
\]
strictly decreasing with $\FAi'(x) = -\mathrm{Ai}(x)^2$.

\medskip
\noindent\textbf{Compactness constants}
(all depend only on $A_0$, not on $c$, $l$):
\[
  c_* := \min_{[-A_0, A_0]} \FAi > 0, \qquad
  c_{\min} := \min_{[-A_0, A_0]} \mathrm{Ai}^2 > 0, \qquad
  C_{\max} := \max_{[-A_0, A_0]} \mathrm{Ai}^2 < \infty.
\]

\medskip
\noindent\textbf{Lipschitz constant:}
$L_{A_0} := \sup_{[-A_0,A_0]} |(\mathrm{Ai}^2)'|
= 2\sup |\mathrm{Ai}\,\mathrm{Ai}'| < \infty$.

\medskip
\noindent\textbf{Airy parameter and connection error.}
Following ORX \cite{OsipovRokhlinXiao}, Ch.~4,
define the Airy parameter
\begin{equation}\label{eq:mul}
  \mu_l := \chi_l / (2c^{4/3}),
\end{equation}
where $\chi_l$ is the $l$-th eigenvalue of the angular
Sturm--Liouville operator.
The ORX Airy-connection error is:
\begin{equation}\label{eq:El}
  E_l(c) := \lambda_l - \FAi(x_l),
\end{equation}
viewed as a function $E_l = E(\mu_l, c)$.

\medskip
\noindent\textbf{Distinction from \cite{PhaseNondeg}.}
The Airy-connection error $E_l(c)$ here (eigenvalue error)
is a different object from the phase error $E_k(c)$
in \cite{PhaseNondeg} (WKB phase error at $s_*$).
Both errors satisfy similar $O(c^{-2/3})$ pointwise bounds
and $O(c^{-4/3})$ discrete stability via the same
ORX regularity argument,
but the objects being approximated are fundamentally
different: eigenvalues vs.\ WKB phases.

% ============================================================
\section{Lemma U: Pointwise Airy Approximation}
\label{sec:lemmaU}
% ============================================================

\begin{lemma}[Pointwise Airy approximation, after ORX]
\label{lem:U}
There exists $C_1 = C_1(A_0) > 0$ such that
for all $c \ge c_0$ and $|l - \NSh| \le A_0 c^{1/3}$:
\begin{equation}\label{eq:U}
  |\lambda_l - \FAi(x_l)| \le C_1\,c^{-2/3}.
\end{equation}
\end{lemma}

\begin{proof}
This is ORX \cite{OsipovRokhlinXiao}, Theorem~4.6,
restated in the notation of this note.
The bound $O(c^{-2/3})$ (rather than $O(c^{-1/3})$)
holds in the eigenvalue approximation because the
eigenvalue $\lambda_l$ is a scalar quantity
(not a phase), and the ORX error estimate for the
concentration eigenvalue is of order $c^{-2/3}$.
\end{proof}

\begin{remark}[Why $c^{-2/3}$ and not $c^{-1/3}$]
In \cite{PhaseNondeg}, the WKB phase error is $O(c^{-1/3})$
because it accumulates over an interval of length $c^{-2/3}$
with an integrand of order $O(1)$.
Here, $\lambda_l$ is the value of the integral itself
(via the Pr\"ufer phase formula), so the error is
one order smaller: $O(c^{-2/3})$.
This is consistent with ORX, Thm.~4.6.
\end{remark}

% ============================================================
\section{Lemma U': Discrete Derivative Stability}
\label{sec:lemmaUprime}
% ============================================================

\begin{lemma}[Discrete Airy derivative stability]
\label{lem:Uprime}
Under the same conditions as Lemma~\ref{lem:U},
there exists $C_2 = C_2(A_0) > 0$ such that
for all $c \ge c_0$ and $|l - \NSh| \le A_0 c^{1/3}$:
\begin{equation}\label{eq:Uprime}
  \bigl|(\lambda_l - \lambda_{l+1})
        - (\FAi(x_l) - \FAi(x_{l+1}))\bigr|
  \le C_2\,c^{-2/3}.
\end{equation}
\end{lemma}

\begin{remark}[Why this does not follow from Lemma~U alone]
Lemma~\ref{lem:U} bounds $|E_l|$ and $|E_{l+1}|$
individually by $C_1 c^{-2/3}$, giving only
$|E_l - E_{l+1}| \le 2C_1 c^{-2/3}$ by triangle inequality.
This is too weak: we need $O(c^{-2/3})$ with an \emph{explicit}
constant not doubling $C_1$, and moreover the
structural argument below gives a much stronger
bound $O(c^{-4/3}) \subset O(c^{-2/3})$, which is what
Paper~VIII requires for the amplitude regularity term
in Corollary~C, Term~2.
\end{remark}

\begin{proof}
Write $E_l(c) = \lambda_l - \FAi(x_l)$ as in~\eqref{eq:El},
so that the left side of~\eqref{eq:Uprime} equals
$|E_l(c) - E_{l+1}(c)|$.

\textbf{Step 1 ($C^1$ regularity in $\mu_l$).}
By ORX \cite{OsipovRokhlinXiao}, Ch.~4, Eq.~(4.31),
the error $E_l = E(\mu_l, c)$ is differentiable in
the Airy parameter $\mu_l = \chi_l/(2c^{4/3})$
with
\[
  \left|\frac{\partial E}{\partial \mu}\right| \le C_A c^{-2/3}
\]
uniformly for $\mu_l$ in the transition window.
(The $c^{-2/3}$ factor reflects that one $\mu$-step
corresponds to one Airy-scale increment of the eigenvalue.)

\textbf{Step 2 (Step size in $\mu_l$).}
The eigenvalue spacing of the Sturm--Liouville operator:
\[
  |\chi_{l+1} - \chi_l| = O(c^{1/3})
\]
(Paper~IV \cite{paperIV}, eigenvalue spacing in the
Galerkin regime; or directly from ORX Ch.~2).
Hence:
\begin{equation}\label{eq:mu-step}
  |\mu_{l+1} - \mu_l|
  = \frac{|\chi_{l+1} - \chi_l|}{2c^{4/3}}
  = \frac{O(c^{1/3})}{2c^{4/3}}
  = O(c^{-1}).
\end{equation}

\textbf{Step 3 (Mean value estimate).}
By Steps 1--2 and the mean value theorem:
\begin{equation}\label{eq:E-discrete-diff}
  |E_{l+1}(c) - E_l(c)|
  \le \left|\frac{\partial E}{\partial \mu}\right|
       \cdot |\mu_{l+1} - \mu_l|
  \le C_A c^{-2/3} \cdot O(c^{-1})
  = O(c^{-5/3}).
\end{equation}
Since $O(c^{-5/3}) = O(c^{-2/3})$,
we can set $C_2 := C_A \cdot C_\mu$ where
$C_\mu$ is the implicit constant in~\eqref{eq:mu-step}.
This gives~\eqref{eq:Uprime}.
\end{proof}

\begin{remark}[Structural origin of the bound]
The bound $O(c^{-5/3})$ in~\eqref{eq:E-discrete-diff}
is actually \emph{better} than the $O(c^{-2/3})$ stated
in~\eqref{eq:Uprime}.
This is the same structural phenomenon seen in
\cite{PhaseNondeg}, Lemma~D: the discrete stability of the
Airy-connection error is $O(c^{-4/3})$ at the phase level,
and here $O(c^{-5/3})$ at the eigenvalue level.
Both are far better than what Paper~VIII requires.
We state only $O(c^{-2/3})$ in~\eqref{eq:Uprime}
because that is the threshold used in Corollary~C
of Paper~VIII.
\end{remark}

% ============================================================
\section{Corollary: ass:gap is Unconditional}
\label{sec:corollary}
% ============================================================

\begin{corollary}[Assumption~3.1 of Paper~VIII]
\label{ass:gap}
There exists $\kappa_0 = \kappa_0(A_0) > 0$ such that
for all $c \ge c_0(A_0)$ and $|l - \NSh| \le A_0 c^{1/3}$:
\[
  \lambda_l - \lambda_{l+1} \ge \kappa_0\,(c/2)^{-1/3}.
\]
\end{corollary}

\begin{proof}
From Lemma~\ref{lem:Uprime} and the MVT applied to $\FAi$:
\[
  \FAi(x_l) - \FAi(x_{l+1})
  = \mathrm{Ai}(\xi_l)^2 \cdot (c/2)^{-1/3},
  \qquad \xi_l \in [x_l, x_{l+1}].
\]
Since $|x_l| \le A_0$ and $\mathrm{Ai}^2 \ge c_{\min} > 0$
on $[-A_0, A_0]$:
\[
  \FAi(x_l) - \FAi(x_{l+1}) \ge c_{\min}\,(c/2)^{-1/3}.
\]
By~\eqref{eq:Uprime}:
\[
  \lambda_l - \lambda_{l+1}
  \ge \FAi(x_l) - \FAi(x_{l+1}) - C_2 c^{-2/3}
  \ge c_{\min}\,(c/2)^{-1/3} - C_2 c^{-2/3}.
\]
For $c \ge c_0$ large enough,
$c_{\min}\,(c/2)^{-1/3} \ge 2 C_2 c^{-2/3}$
(since $(c/2)^{-1/3}/c^{-2/3} = (c/2)^{-1/3} c^{2/3}
= 2^{1/3} c^{1/3} \to \infty$).
Set $\kappa_0 := c_{\min}/2$.
\end{proof}

\begin{remark}[What this resolves in Paper~VIII]
With Corollary~\ref{ass:gap}:
\begin{itemize}
  \item Assumption~3.1 (ass:gap) of Paper~VIII \cite{paperVIII}
        is proved unconditionally.
  \item Lemma~F (Dyadic Separation) holds unconditionally.
  \item Corollary~C(a) holds unconditionally.
  \item The remaining open problems in Paper~VIII reduce to:
        \textbf{B-strong} and \textbf{Proposition~U(U')}
        as a \emph{proof} (vs.\ the reference to this note).
\end{itemize}
\end{remark}

\begin{remark}[Precise dependency chain]
The full chain is:
\[
  \text{ORX Thm.~4.6}
  \;\xrightarrow{\text{pointwise}}\;
  \text{Lemma~U}
  \;\xrightarrow{\partial_\mu \text{ regularity}+\text{step size}}
  \;\text{Lemma~U'}
  \;\xrightarrow{c_{\min}>0}\;
  \text{ass:gap}.
\]
The \cite{PhaseNondeg} note is \emph{not} in this chain.
It enters Paper~VIII only via Corollary~2
(H1 unconditional), which is logically independent
of the eigenvalue gap structure.
\end{remark}

% ============================================================
\begin{thebibliography}{9}

\bibitem{paperIV}
  [Author],
  \textit{Semiclassical Equidistribution of PSWF Densities
  via Pr\"ufer Phase Analysis},
  Paper~IV, prolate-gram-coercivity, 2026.

\bibitem{paperVIII}
  [Author],
  \textit{Scale-Separated Dyadic Cancellation for PSWF
  Galerkin Operators},
  Paper~VIII (working draft), prolate-gram-coercivity, 2026.

\bibitem{PhaseNondeg}
  [Author],
  \textit{Phase Non-Degeneracy of the PSWF WKB Phases:
  $\alpha^{(c)} = \pi/2 + O(c^{-1/3})$},
  companion note (\texttt{phase\_nondeg\_lemma.tex}),
  prolate-gram-coercivity, 2026.

\bibitem{OsipovRokhlinXiao}
  A.~Osipov, V.~Rokhlin, H.~Xiao,
  \textit{Prolate Spheroidal Wave Functions of Order Zero},
  Springer, 2013.

\end{thebibliography}

\end{document}
